% ============================================================================
% Chapter 17: Hardware in Detail
% ============================================================================

\chapter{Hardware in Detail}
\label{ch:hardware-detail}

\begin{objectives}
  \item Understand the internal components of the CPU
  \item Explain the difference between RAM and ROM
  \item Describe how the ALU and Control Unit work together
  \item Compare different types of memory and storage
\end{objectives}

\bytetip[tip]{Time to look \emph{inside} the computer! We're going to open up the CPU and see what's really going on in there. Don't worry --- no screwdrivers needed!}

\section{Inside the CPU}
\label{sec:inside-cpu}
\index{CPU!components}\index{ALU}\index{Control Unit}

The CPU has two main parts that work together:

% --- Figure 17.1: Inside the CPU Block Diagram ---
\begin{figure}[H]
\centering
\begin{tikzpicture}
  % CPU outer box
  \draw[NavyBlue, line width=2.5pt, rounded corners=8pt] (-0.5,-1.5) rectangle (10.5,4.5);
  \node[font=\large\bfseries\sffamily, text=NavyBlue, anchor=north west] at (-0.2,4.3) {CPU};
  
  % ALU box
  \node[draw=FunCoral, fill=FunCoral!10, rounded corners=6pt,
        minimum width=4cm, minimum height=2.5cm, line width=1.5pt,
        font=\small\sffamily, align=center] (alu) at (2.5,2) {
    {\Large\faCalculator}\\[3pt]
    \textbf{ALU}\\
    {\scriptsize Arithmetic Logic Unit}\\[2pt]
    {\tiny Does all maths \& logic}
  };
  
  % Control Unit box
  \node[draw=FunSky, fill=FunSky!10, rounded corners=6pt,
        minimum width=4cm, minimum height=2.5cm, line width=1.5pt,
        font=\small\sffamily, align=center] (cu) at (7.5,2) {
    {\Large\faDirections}\\[3pt]
    \textbf{Control Unit (CU)}\\[2pt]
    {\tiny Directs all operations}
  };
  
  % Registers
  \node[draw=FunPurple, fill=FunPurple!10, rounded corners=4pt,
        minimum width=8cm, minimum height=0.8cm, line width=1pt,
        font=\small\sffamily] (reg) at (5,0) {
    \faMemory\ \textbf{Registers} --- Ultra-fast tiny storage
  };
  
  % RAM outside
  \node[hwbox, fill=FunMint!10, draw=FunMint, minimum width=3cm] (ram) at (5,-3) {
    \faDatabase\ \textbf{RAM}\\External Memory
  };
  
  % Arrows between components
  \draw[FunPurple, line width=1.5pt, <->] (alu.south) -- ([yshift=0.1cm]reg.north -| alu);
  \draw[FunPurple, line width=1.5pt, <->] (cu.south) -- ([yshift=0.1cm]reg.north -| cu);
  \draw[NoteOrange, line width=1.5pt, <->] (alu.east) -- (cu.west)
    node[midway, above, font=\tiny\sffamily\bfseries, text=NoteOrange] {Data \& Control};
  \draw[FunMint, line width=2pt, <->] (reg.south) -- (ram.north)
    node[midway, right, font=\tiny\sffamily, text=FunMint] {Data transfer};
\end{tikzpicture}
\caption{Inside the CPU: ALU, Control Unit, and Registers}
\label{fig:inside-cpu}
\end{figure}

\section{RAM vs ROM}
\label{sec:ram-rom}
\index{RAM}\index{ROM}

% --- Figure 17.2: RAM vs ROM Comparison ---
\begin{figure}[H]
\centering
\begin{tikzpicture}
  % RAM Panel
  \fill[FunSky!12] (0,0) rectangle (5.5,5.5);
  \draw[FunSky, line width=2pt] (0,0) rectangle (5.5,5.5);
  \node[font=\large\bfseries\sffamily, text=FunSky] at (2.75,5) {\faDatabase\ RAM};
  \node[font=\small\sffamily, text=NavyBlue, align=left, anchor=north west] at (0.3,4.3) {
    \textcolor{FunSky}{\faAngleRight}\ Volatile\\[3pt]
    \textcolor{FunSky}{\faAngleRight}\ Data lost when off\\[3pt]
    \textcolor{FunSky}{\faAngleRight}\ Temporary storage\\[3pt]
    \textcolor{FunSky}{\faAngleRight}\ Fast read/write\\[3pt]
    \textcolor{FunSky}{\faAngleRight}\ Example: 8 GB RAM
  };
  
  % ROM Panel
  \fill[FunPurple!10] (6.5,0) rectangle (12,5.5);
  \draw[FunPurple, line width=2pt] (6.5,0) rectangle (12,5.5);
  \node[font=\large\bfseries\sffamily, text=FunPurple] at (9.25,5) {\faLock\ ROM};
  \node[font=\small\sffamily, text=NavyBlue, align=left, anchor=north west] at (6.8,4.3) {
    \textcolor{FunPurple}{\faAngleRight}\ Non-volatile\\[3pt]
    \textcolor{FunPurple}{\faAngleRight}\ Data stays forever\\[3pt]
    \textcolor{FunPurple}{\faAngleRight}\ Permanent instructions\\[3pt]
    \textcolor{FunPurple}{\faAngleRight}\ Read-only (mostly)\\[3pt]
    \textcolor{FunPurple}{\faAngleRight}\ Example: BIOS chip
  };
  
  % Center overlap
  \node[draw=NavyBlue, fill=white, rounded corners=4pt, 
        font=\scriptsize\sffamily, text=NavyBlue, align=center,
        inner sep=4pt] at (6, 2.75) {Both are\\memory\\inside the\\computer};
\end{tikzpicture}
\caption{RAM vs ROM --- two types of memory with very different jobs}
\label{fig:ram-vs-rom}
\end{figure}

\begin{theorybox}[RAM vs ROM Quick Summary]
\begin{tabular}{lll}
  & \textbf{RAM} & \textbf{ROM} \\
  \midrule
  Full name & Random Access Memory & Read-Only Memory \\
  Volatile? & Yes (loses data) & No (keeps data) \\
  Speed & Very fast & Slower than RAM \\
  Used for & Running programs & Boot instructions \\
\end{tabular}
\end{theorybox}

\bytetip[tip]{Here's an easy way to remember: \textbf{RAM} is like your desk --- it holds what you're currently working on. \textbf{ROM} is like a recipe book --- the information is permanent and doesn't change.}

\begin{funfact}
The first computers had only a few \textbf{kilobytes} of RAM. Today, even budget smartphones have 4--6 GB of RAM --- that's over a million times more! Modern gaming PCs can have 64 GB or even 128 GB of RAM.
\end{funfact}

\begin{reviewqs}
  \item What are the two main parts inside the CPU?
  \item What does ALU stand for and what does it do?
  \item What is the difference between RAM and ROM?
  \item What are registers in a CPU?
  \item Why is RAM called ``volatile'' memory?
\end{reviewqs}

\begin{challenge}
Check your school computer's RAM: press \key{Windows}+\key{I} $\rightarrow$ System $\rightarrow$ About. How much RAM does it have? Research online: how much RAM does a typical smartphone have? Compare the two and write three sentences about your findings.
\end{challenge}
